\section{Switching and cycle holonomy}
\label{sec:holonomy}

Assume throughout this section that \(P=|H|\) and \(G_P\) is connected.

\begin{theorem}[Holonomy classification]\label{thm:holonomy}
For a fixed \(\sigma\in\{\pm1\}\), the following are equivalent:
\begin{enumerate}[label=(\roman*),leftmargin=2.4em]
\item \(H=\sigma DPD^*\) for a diagonal unitary \(D\);
\item there are phases \(d_i\) such that
\[
 \phi_{ij}=\sigma d_i\overline{d_j}
 \quad\text{on every edge};
\]
\item every oriented cycle \(C\) satisfies
\begin{equation}\label{eq:cycle-condition}
 \Hol_\phi(C)=\sigma^{|C|},
 \quad\text{equivalently}\quad
 \Hol_\sigma(C)=1.
\end{equation}
\end{enumerate}
\end{theorem}

\begin{proof}
The first two statements are entrywise equivalent.  Their phases
telescopically give \eqref{eq:cycle-condition}.

Conversely, fix a root and a spanning tree.  Set the root phase
arbitrarily and propagate \(d_i\) along the unique tree paths using
\(\phi_{ij}=\sigma d_i\overline{d_j}\).  For a chord, its fundamental
cycle condition says that the already assigned endpoint phases satisfy
the same equation.  Hence the formula holds on every edge.
\end{proof}

\begin{definition}[Balanced and antibalanced sectors]
The phase pattern is \emph{balanced} when the theorem holds for
\(\sigma=+1\), and \emph{antibalanced} when it holds for
\(\sigma=-1\).
\end{definition}

For a cycle of length \(\ell\), balanced means holonomy \(1\), whereas
antibalanced means holonomy \((-1)^\ell\).  On an odd graph the two
conditions are distinct; on a bipartite graph they coincide after a
bipartition sign switch.

\begin{corollary}[Fundamental-cycle census]\label{cor:fundamental}
Let \(T\) be any spanning tree of \(G_P\).  It is enough to test
\eqref{eq:cycle-condition} on the
\[
 |E|-|V|+1
\]
fundamental cycles formed by the chords \(e\notin T\).
\end{corollary}

\begin{proof}
The spanning-tree construction in \cref{thm:holonomy} uses exactly
those conditions.  All other cycle holonomies are products of
fundamental ones and their conjugates.
\end{proof}

\begin{remark}[Gauge normal form]
A diagonal unitary can switch every tree-edge phase to \(+1\).  The
remaining chord phase is, up to orientation conjugation, the holonomy
of its fundamental cycle.  Thus cycle rank, not edge count, is the
number of independent scalar phase defects.
\end{remark}
